\chapter{课程总览与全书导航}
\label{chap:course_overview}

本书来自 UC Berkeley / Berkeley RDI 课程《CS294/194-280: Advanced Large Language Model Agents, Spring 2025》的 12 次公开 lecture。它不是把视频逐页转写后堆在一起，而是试图把整门课组织成一套围绕 \textbf{advanced LLM agents} 的教材化知识结构。

全书的核心主张可以概括为一句话：\textbf{高级 LLM agent 的能力，来自模型、推理时计算、外部工具、环境反馈、记忆结构、形式验证与安全边界之间的系统协同，而不只是更大的参数规模。}

围绕这条主线，本书分成六个 Parts：

\begin{itemize}
\item \textbf{Part I: Foundations of Advanced LLM Agents} 讨论推理时计算、学习推理、以及语言智能体的记忆与规划，建立全书的基础词汇表。
\item \textbf{Part II: Agentic Workflows, Tools, and Code} 把推理技术放回后训练 recipes、coding agents 与漏洞检测这些真实工作流里。
\item \textbf{Part III: Multimodal and Interactive Agents} 把讨论从纯文本扩展到 web、GUI、OS 与多模态环境中的 perception-action loop。
\item \textbf{Part IV: Formal Mathematics, Verification, and Theorem Proving} 展示最强的 environment-grounded reasoning 形态，即由 proof assistant 和 formal system 提供机械可检验反馈的智能体系统。
\item \textbf{Part V: Abstraction, Discovery, Safety, and Security} 把问题进一步推向抽象发现、科学探索以及 agentic AI 的安全与权限边界。
\item \textbf{Part VI: Cross-Course Extensions and 2025 Updates} 把 Berkeley Fall 2024、Berkeley Fall 2025 与 Stanford CS329A 的官方课程材料并入一个扩展部分，用来回答“如果读完 Spring 2025 主线之后，还要继续补系统、评测、自改进与最新公开课程线索，应当往哪里走”。
\end{itemize}

从阅读路径看，若你第一次系统接触这门课，建议按顺序阅读；若你更关心某个主题，也可以直接沿以下线索切入：

\begin{itemize}
\item 推理与搜索主线：第 1--4 章，再跳到第 9--11 章。
\item 实际 agent workflow 主线：第 4--8 章，再看第 12--13 章的风险与治理。
\item 数学与形式化主线：第 9--11 章，再回看第 1--3 章理解它们为何依赖 inference-time search 和 verifier。
\item 安全与系统边界主线：第 6、8、12、13 章。
\item 课程扩展与最新材料主线：在读完第 13 章后继续看第 14--16 章，把 Berkeley 2024/2025 与 Stanford 2025 的课程更新接回本书主线。
\end{itemize}

阅读本书时可以同时参考 lecture-level sidecars。每章背后都保留了 \texttt{source\_manifest.json}、\texttt{coverage\_units.jsonl}、\texttt{figure\_manifest.json}、\texttt{omission\_log.jsonl}、\texttt{eval\_report.json} 等结构化记录，因此你不仅能读到中文教材正文，还能回查它如何从视频、slides 和 readings 生成而来。

最后需要强调的是：本书中的“agent”不是单一算法名词，而是一个系统工程对象。不同章节分别从 reasoning、training、tool use、multimodality、formal verification 和 security 的角度切入；扩展部分则进一步把它放回 Berkeley 与 Stanford 的相邻课程生态中。只有把这些角度放在一起，advanced LLM agents 才真正成为一个完整研究领域。
